% Standalone problem document, generated from the adjacent README.md.
% Compile with XeLaTeX. Mathematical content is maintained in Markdown.
\documentclass[11pt,a4paper]{article}
\usepackage[fontsize=10.5pt]{scrextend}
\usepackage[margin=22mm,headheight=15pt,headsep=9mm,footskip=11mm]{geometry}
\usepackage{fontspec}
\setmainfont{texgyrepagella}[Extension=.otf,UprightFont=*-regular,BoldFont=*-bold,ItalicFont=*-italic,BoldItalicFont=*-bolditalic]
\setsansfont{texgyreheros}[Extension=.otf,UprightFont=*-regular,BoldFont=*-bold,ItalicFont=*-italic,BoldItalicFont=*-bolditalic]
\setmonofont{texgyrecursor}[Extension=.otf,UprightFont=*-regular,BoldFont=*-bold,ItalicFont=*-italic,BoldItalicFont=*-bolditalic,Scale=MatchLowercase]
\usepackage{amsmath,amssymb}
\usepackage{unicode-math}
\setmathfont{texgyrepagella-math.otf}
\usepackage{xcolor,microtype,enumitem,needspace,fancyhdr}
\usepackage{bookmark}
\definecolor{ink}{HTML}{17344B}
\definecolor{accent}{HTML}{197D80}
\definecolor{muted}{HTML}{596773}
\hypersetup{colorlinks=true,linkcolor=accent,urlcolor=accent,pdftitle={TR-25 - Cohen--Macaulay
coordinate rings for every tensor border-rank
variety},pdfauthor={Open Problems in Numerical Linear Algebra}}
\urlstyle{same}
\setlength{\parindent}{0pt}
\setlength{\parskip}{5pt plus 1pt minus 1pt}
\linespread{1.04}
\setlength{\emergencystretch}{3em}
\widowpenalty=10000
\clubpenalty=10000
\displaywidowpenalty=10000
\allowdisplaybreaks[1]
\setlist{nosep,leftmargin=1.5em}
\providecommand{\tightlist}{\setlength{\itemsep}{0pt}\setlength{\parskip}{0pt}}
\providecommand{\pandocbounded}[1]{#1}
\pagestyle{fancy}
\fancyhf{}
\fancyhead[L]{\sffamily\footnotesize\color{muted}OPEN PROBLEMS IN NUMERICAL LINEAR ALGEBRA}
\fancyhead[R]{\sffamily\bfseries\footnotesize\color{accent}TR-25}
\fancyfoot[L]{\parbox[b]{0.9\textwidth}{\sffamily\fontsize{8}{10}\selectfont\color{muted}Editorial ratings; a bounded search cannot certify that a problem remains unsolved.\\Literature check: 2026-09-10}}
\fancyfoot[R]{\sffamily\footnotesize\color{muted}\thepage}
\renewcommand{\headrulewidth}{0.3pt}
\renewcommand{\headrule}{\hbox to\headwidth{\color{accent}\leaders\hrule height \headrulewidth\hfill}}
\makeatletter
\renewcommand{\section}{\Needspace{5\baselineskip}\@startsection{section}{1}{0pt}{12pt plus 2pt minus 2pt}{4pt}{\sffamily\bfseries\large\color{ink}}}
\renewcommand{\subsection}{\Needspace{4\baselineskip}\@startsection{subsection}{2}{0pt}{9pt plus 1pt minus 1pt}{3pt}{\sffamily\bfseries\normalsize\color{ink}}}
\makeatother
\setcounter{secnumdepth}{0}
\begin{document}
{\sffamily\small\color{muted}Tensor computations\par}
\vspace{3pt}
{\raggedright\hyphenpenalty=10000\exhyphenpenalty=10000\sffamily\bfseries\fontsize{21}{24}\selectfont\color{ink}Cohen--Macaulay
coordinate rings for every tensor border-rank variety\par}
\vspace{7pt}
{\sffamily\small\textbf{Difficulty:} extreme\quad\textbf{Importance:} interesting
to specialist\par}
{\sffamily\small\bfseries\color{accent}Status: Partially resolved\par}
\vspace{4pt}
{\color{accent}\hrule height 0.8pt}
\vspace{6pt}
\textbf{Rating rationale:} An all-format theorem would overcome
longstanding obstructions in the local algebra of secant varieties and
encompass other major Cohen--Macaulay conjectures. The depth and
resolution questions chiefly concern specialists in tensor algebraic
geometry.

\subsection{Problem statement}\label{problem-statement}

Fix integers \(k\geq3\), \(n_1,\ldots,n_k\geq1\), and \(r\geq1\). In
\(V=\mathbb C^{n_1}\otimes\cdots\otimes\mathbb C^{n_k}\), let \[
X_r=\overline{\left\{
\sum_{\ell=1}^{r}v_{1,\ell}\otimes\cdots\otimes v_{k,\ell}:
v_{i,\ell}\in\mathbb C^{n_i}
\right\}}^{\,\mathrm{Zar}}.
\] Let \(S=\mathbb C[V]\) be the polynomial ring in the tensor entries,
let \(I(X_r)\) be its vanishing ideal, and let \(A_r=S/I(X_r)\).

Is \(A_r\) Cohen--Macaulay for every such format and rank threshold?

Explicitly, if \(\mathfrak m\) is the ideal of positive-degree elements
of \(A_r\), the requested equality is \[
\operatorname{depth}(A_r)_{\mathfrak m}
=\dim(A_r)_{\mathfrak m}.
\] Depth is the maximum length of a regular sequence in the maximal
ideal: each successive element is a non-zero-divisor after quotienting
by its predecessors. Dimension is Krull dimension. This is Oeding's
conjecture that all secant varieties of Segre products are
arithmetically Cohen--Macaulay. The zero tensor is retained by using the
affine cone.

\subsection{Why it matters}\label{why-it-matters}

These are the algebraic sets underlying low-border-rank tensor
approximation. Cohen--Macaulayness constrains their equations and free
resolutions, supporting the determination of generators and the analysis
of degeneracies. The conjecture does not assert that exact low-rank
tensor approximation is computationally easy or always attained.

\subsection{References and status
check}\label{references-and-status-check}

\begin{enumerate}
\def\labelenumi{\arabic{enumi}.}
\tightlist
\item
  L. Oeding, \emph{Are all Secant Varieties of Segre Products
  Arithmetically Cohen-Macaulay?},
  \href{https://arxiv.org/pdf/1603.08980v2}{Preprint (2016), v2},
  Conjecture 1.1, p.1; §2 collects established cases.
\item
  J. Jagiełła and J. Jelisiejew, \emph{Unrestrictions and concise secant
  varieties}, \href{https://arxiv.org/html/2604.24879v2}{2026 preprint,
  v2}, Conjecture 3.32 and its surrounding discussion; §1.5.3 explains a
  connection to the Salmon ideal problem.
\end{enumerate}

\subsubsection{Status check ---
2026-09-10}\label{status-check-2026-09-10}

Checked Oeding v2, including Theorems 3.1--3.2 and Proposition 3.5, and
Jagiełła--Jelisiejew v2, Conjecture 3.32, with targeted later-resolution
searches. Oeding proves substantive tensor cases, including the fourth
secant of \(\mathbb P^2\times\mathbb P^2\times\mathbb P^n\) for
\(n\ge3\), rather than only the matrix or rank-one cases. The June 2026
source still treats the all-format conjecture as open. Inheritance
results require extra hypotheses and do not establish it universally. No
full proof or counterexample was located.
\end{document}
